\section{TPC--BTOF association QA}
\label{sec:tofqa}

\paragraph{Match flag and $\beta$.} Every QA track with a \fcmd{BTofPidTraits} entry has
\fcmd{btofMatchFlag}$>0$ (\nFlagZero{} entries with flag 0), so the picoDst keeps only matched tracks.
\nFlagOne{} (\fracFlagOne\%) are one-to-one matches and \nFlagTwo{} (\fracFlagTwo\%) have flag 2,
i.e.\ the hit had more than one track candidate (Fig.~\ref{fig:tof_matchflag}).
\nBetaLost{} matched tracks (\fracBetaLost\%) have $\beta=0$ (no valid time) and are dropped from the
TOF-matched sample. The $\beta$ distribution (Fig.~\ref{fig:tof_beta}) peaks just below 1 with a tail
above 1 containing \fracBetaGtOne\% of the matched tracks; these give $m^2<0$.

\paragraph{Local coordinates.} $z_{\rm local}$ is flat with sharp edges at $\pm3.05$\,cm, the matching window
of \fcmd{StBTofMatchMaker} (Fig.~\ref{fig:tof_zlocal}). $y_{\rm local}$ (Fig.~\ref{fig:tof_ylocal}) has a plateau
over the pad, $|y_{\rm local}|\lesssim1.6$\,cm, falling over about 1\,cm on each side, and flat tails out to
$\pm5$\,cm about a factor 50 below the plateau (read by eye from the figure); \fracYLocalIn\% of the entries are within
$|y_{\rm local}|<1.8$\,cm, RMS \rmsYLocal\,cm. The tails broaden at low momentum, where multiple
scattering and the projection uncertainty are largest (Fig.~\ref{fig:tof_ylocal_p}).
\caveat{The flat $y_{\rm local}$ tails are the region where random matches are expected to dominate.
Their size relative to the plateau is a direct handle for the mixed-event normalisation.}

\begin{figure}[H]
\centering
\figpanel{0.32}{tof_matchflag}{\fcmd{btofMatchFlag} of QA tracks with BTofPidTraits.}
\hfill
\figpanel{0.32}{tof_beta}{$\beta$ of matched tracks, log scale; red dashed line: $\beta=1$. The
$\beta=0$ entries are in the first bin.}
\hfill
\figpanel{0.32}{tof_ylz}{$y_{\rm local}$ vs $z_{\rm local}$, log colour scale.}
\\[1ex]
\figpanel{0.32}{tof_ylocal}{$y_{\rm local}$ of matched tracks, 0.05\,cm bins, log scale.}
\hfill
\figpanel{0.32}{tof_zlocal}{$z_{\rm local}$ of matched tracks, 0.05\,cm bins, log scale.}
\hfill
\figpanel{0.32}{tof_ylocal_p}{$y_{\rm local}$ vs $p$, log colour scale.}
\\[1ex]
\figpanel{0.32}{tof_zlocal_p}{$z_{\rm local}$ vs $p$, log colour scale.}
\caption{Match flag, $\beta$ and local hit coordinates for QA tracks matched to BTOF (flag$>0$).}
\end{figure}

\paragraph{Particle identification.} $1/\beta$ vs $p$ (Fig.~\ref{fig:tof_invbeta}) shows clean $\pi$, K and p bands
that follow the expected curves, plus a deuteron band and a uniform background between bands.
Fitting Gaussians to the \mtwo{} projection for $0.4<p<1.0$\,\GeVc{} (Fig.~\ref{fig:tof_m2_proj}) gives
\begin{center}
\begin{tabular}{lccc}
\toprule
 & mean (GeV$^2/c^4$) & $\sigma$ (GeV$^2/c^4$) & PDG $m^2$ \\
\midrule
$\pi$ & \mTwoMeanPi & \mTwoSigmaPi & \mTwoPdgPi \\
K & \mTwoMeanK & \mTwoSigmaK & \mTwoPdgK \\
p & \mTwoMeanP & \mTwoSigmaP & \mTwoPdgP \\
d & \mTwoMeanD & \mTwoSigmaD & \mTwoPdgD \\
\bottomrule
\end{tabular}
\end{center}
The $\pi$, K and p positions agree with PDG values within 0.005\,GeV$^2/c^4$, i.e.\ the TOF
calibration in P24iy is in good shape for this run. The deuteron fit has low statistics in this
momentum window and a broad, asymmetric peak; it is shown for orientation only.

\begin{figure}[H]
\centering
\figpanel{0.48}{tof_invbeta}{$1/\beta$ vs $p$; red dashed curves are the expectations for $\pi$, K, p.}
\hfill
\figpanel{0.48}{tof_invbeta_pq}{$1/\beta$ vs $p\times q$.}
\\[1ex]
\figpanel{0.48}{tof_m2_p}{\mtwo{} vs $p$.}
\hfill
\figpanel{0.48}{tof_m2_pq}{\mtwo{} vs $p\times q$; d appears only for positive charge.}
\\[1ex]
\figpanel{0.48}{tof_m2_proj}{\mtwo{} for $0.4<p<1.0$\,\GeVc{} (black points), with Gaussian fits around the
$\pi$, K, p and d peaks (red). Log scale; 0.01\,GeV$^2/c^4$ bins.}
\hfill
\figpanel{0.48}{tof_dinvbeta_pi}{$1/\beta-1/\beta_{\rm exp}^{\pi}$ vs $p$.}
\caption{TOF PID for the TOF-matched sample (flag$>0$, $\beta>0$), log colour scale.}
\end{figure}

\begin{figure}[H]
\centering
\figpanel{0.48}{tof_nsig_pi}{$n\sigma^{\rm TOF}_\pi$ vs $p$.}
\hfill
\figpanel{0.48}{tof_nsig_k}{$n\sigma^{\rm TOF}_{\rm K}$ vs $p$.}
\\[1ex]
\figpanel{0.48}{tof_nsig_p}{$n\sigma^{\rm TOF}_{\rm p}$ vs $p$.}
\caption{Approximate TOF $n\sigma$ with a fixed $1/\beta$ resolution of 0.013 (Sec.~\ref{sec:tofdef}); not calibrated.}
\end{figure}

\paragraph{Occupancy.} Figure~\ref{fig:tof_occ_hits} shows the tray--module occupancy of all stored hits.
The median module has \medianModHits{} hits in this sample. Trays \deadTrayList{} are empty;
trays with empty modules, given as tray(number of empty modules), are \partTrayList.
In total \nDeadModule{} of 3840 modules (\fracDeadModule\%) are empty, and no matched track points to an
empty module (\nMatchedInEmpty{} entries), as expected. In each of the four trays with 8 empty modules, the empty modules are \eightModPattern{} (same list in all four). Isolated low-occupancy modules are also visible.
The global hit-position map of matched tracks (Fig.~\ref{fig:tof_hitpos}) shows the same gaps as
horizontal white bands at fixed $\phi$, separately for $z<0$ (east) and $z>0$ (west).
Cell occupancy (1--6) is uniform within a factor \cellMaxOverMin{} (max/min).

\caveat{The dead-tray/module list is from one run only. A run-by-run (or time-period) BTOF status map is
needed for the mixed-event step, since mixing events with different dead areas biases the
random-match rate.}

\begin{figure}[H]
\centering
\figpanel{0.48}{tof_occ_hits}{Tray vs module for all stored BTOF hits (linear colour scale).}
\hfill
\figpanel{0.48}{tof_occ_matched}{Tray vs module of the cell matched to a QA track (flag$>0$), linear colour scale.}
\\[1ex]
\figpanel{0.48}{tof_hitpos}{Global $\phi$ vs $z$ of the BTOF hit position of matched tracks, log colour scale.}
\caption{BTOF occupancy.}
\end{figure}

\paragraph{Matching efficiency.} Table~\ref{tab:eff} and Fig.~\ref{fig:tof_eff_pt} give
$\varepsilon_{\rm match}$ versus $p_T$. It rises from \effLow{} at $0.15<p_T<0.3$\,\GeVc{} to a plateau of
0.63--0.66 above 0.5\,\GeVc{} (Table~\ref{tab:eff}). Versus $\eta$ (Fig.~\ref{fig:tof_eff_eta}) it is \effEtaCentral{} for $|\eta|<0.5$
with a shallow dip at $\eta\approx0$ and maxima near $|\eta|\approx0.5$, then drops to \effEtaOuter{} averaged
over $0.5<|\eta|<1$ as the edge of the BTOF acceptance is reached. Versus $\phi$ (6$^\circ$ bins) it ranges
from \effPhiMin{} (at $\phi\approx\effPhiMinAt$) to \effPhiMax.

\begin{table}[h]
\centering
\caption{TOF matching efficiency vs $p_T$ for QA tracks (generated table \fpath{tables/eff.tex}). Statistical errors.}
\label{tab:eff}
\input{tables/eff.tex}
\end{table}

\begin{figure}[H]
\centering
\figpanel{0.32}{tof_eff_pt}{$\varepsilon_{\rm match}$ vs $p_T$ (0.05\,\GeVc{} bins).}
\hfill
\figpanel{0.32}{tof_eff_eta}{$\varepsilon_{\rm match}$ vs $\eta$.}
\hfill
\figpanel{0.32}{tof_eff_phi}{$\varepsilon_{\rm match}$ vs $\phi$ (1$^\circ$ bins).}
\caption{TOF matching efficiency $N_{\rm TOF}/N_{\rm TPC}$ (flag$>0$, $\beta>0$ over QA tracks). Error bars:
68\% Bayesian binomial intervals.}
\end{figure}
